\section{周期轨道设计}
\label{sec:orbit}

\subsection{轨道族总览}

地月空间的可用轨道按离月距离分为四层，尺度是主线：离月越远，地球引力
占比越高，动力学从二体开普勒过渡到三体问题，轨道从稳定变为不稳定
（表~\ref{tab:families}）。e2m2e 提供统一的轨道族生成入口
\texttt{orbit\_family\_generation()}，覆盖八类轨道族；Lyapunov 与共振
轨道（RO）在微分修正策略层额外支持。

\begin{table}[htbp]
\centering\small
\caption{周期轨道族一览}
\label{tab:families}
\begin{tabularx}{\textwidth}{@{}llX@{}}
\toprule
\textbf{轨道族} & \textbf{位置} & \textbf{特点与用途} \\
\midrule
DRO & 月心 & 远距逆行，中性稳定、维持成本低，月轨空间站候选 \\
DPO & L2 附近 & 顺行远距，不稳定，可用于月球背面通信中继 \\
Halo / NRHO & L1/L2 & 三维周期；NRHO 为高振幅近直线形态（$A_z > 0.1$ 无量纲），
Gateway 采用 \\
Axial & L1/L2 & Gómez Type B 分岔，轴向运动，覆盖特性与 Halo 互补 \\
Lyapunov & L1/L2/L3 & 平面周期轨道 \\
Lissajous & L1/L2/L3 & 拟周期（面内/面外频率不可通约），非线性中心
约化流上生成 \\
SPO / LPO & L4/L5 & 短周期 / 长周期轨道，大振幅 LPO 呈马蹄形 \\
Horseshoe & L4/L5 & LPO 链的大振幅成员 \\
RO & --- & 共振轨道，$y$ 轴对称出发 \\
\bottomrule
\end{tabularx}
\end{table}

族的周期语义严格区分：Halo、NRHO、Axial、SPO、LPO、Horseshoe、DRO
成员严格周期；Lissajous 成员为拟周期有界轨迹，不做周期闭合。

\subsection{初值猜测：从解析近似出发}

周期轨道设计的第一块基石是初猜。e2m2e 的初猜来源有三：

\begin{itemize}
  \item \textbf{Richardson 三阶解析近似}（1980）：共线平动点附近 Halo
        轨道的三阶解析解。内部流程为：求解平动点距离参数 $\gamma$
        （五次方程，Brent 法）$\to$ 计算面内频率 $\omega_p$ $\to$
        组装三阶系数与解析解。小振幅种子（如无量纲 $z_0 = 0.001$）
        精度最高，交由延拓放大；
  \item \textbf{种子轨道}：DRO 从标准种子（振幅约 90786 km）出发，
        固定 $x$ 轴穿越点做微分修正并双向延拓，包络覆盖
        1737--110000 km；
  \item \textbf{Normal Form 表征参数}：把 CR3BP 平动点附近的动力学
        化简为作用量--角变量形式的表征参数
        $[q_1, p_1, I_2, \theta_2, I_3, \theta_3]$，用几个常数描述
        本需六个状态量随时间演化的轨道。流水线为：星历轨道初值 $\to$
        动力学替代轨道 $\to$ quasi-Floquet 变换 $\to$ 中心流形化简
        $\to$ 表征参数；其中 quasi-Floquet 化简可选 36 维矩阵法或
        21 维 $\mathfrak{sp}(6)$ Lie 代数参数化（构造上保辛），中心
        流形 Lie 变换默认截断 10 阶。
\end{itemize}

\subsection{微分修正：利用对称性的 STM 牛顿迭代}

微分修正迭代调整初始条件使轨道满足周期性约束，核心是状态转移矩阵
（STM）牛顿迭代。e2m2e 充分利用 CR3BP 的三种镜像对称性来缩小问题
规模——例如 XZ 平面对称意味着：若
$(x, y, z, \dot x, \dot y, \dot z)$ 是解，则
$(x, -y, z, -\dot x, \dot y, -\dot z)$ 也是解；于是 Halo 轨道只需
从 XZ 平面（$y = 0$）垂直出发，修正到半周期后再次垂直穿越即可。

对称性与固定参数的不同组合构成 11 种修正策略（Lyapunov/DRO 固定
$x_0$、固定周期轨道固定 $T_{1/2}$、RO 固定 $y_0$、Halo 固定 $z_0$ 或
$x_0$、Axial 固定 $v_{z0}$、L4/L5 的 SPO/LPO 无对称全周期闭合等）。
架构上策略层与求解器严格分离：策略函数只生成不可变配置对象，执行
下沉到 Rust CR3BP 修正内核。

收敛失败可诊断：历史记录给出每步误差、终止原因（正常收敛 / 修正量低于
机器精度 / 发散 / 雅可比奇异 / 停滞 / 收敛到 $T \approx 0$ 的寄生根），
雅可比奇异提示约束线性相关、需要更换策略。

\subsection{多重打靶：长弧段的星历修正}

微分修正对初猜敏感，弧段越长越甚。多重打靶把轨道分为 $N$ 段独立积分，
在拼接点施加连续性约束
\begin{equation}
  \vect{x}_i(t_{i+1}) = \vect{x}_{i+1}(t_{i+1}),
\end{equation}
通过同时调整各段初始状态（可选含时间节点）使所有约束同时满足，把
“长弧敏感”转化为“多段小修正”。标准实现的容差达 $10^{-10}$。

星历模型下的 patch 点修正统一走 Rust 多重打靶内核，并按轨道稳定性
分派两条路径：

\begin{itemize}
  \item \textbf{稳定轨道}（DRO）：速度加权的两级路径——修正一圈后自由
        外推有界；
  \item \textbf{不稳定轨道}（Halo/NRHO/DPO）：全程分段打靶，逐段积分
        填满星历网格；NRHO 每段固定 1 圈，Halo 最长 3 圈一段，DPO
        每圈 64 个等时节点。
\end{itemize}

节点采样策略同样按轨道物理定制：Halo 采用近月点加密（近月窗口内
5 节点 + 窗外 8 节点），因为近月点速度大、STM 病态；NRHO 等时间采样。

\subsection{延拓：从一条轨道到一个族}

给定一条已收敛的种子轨道，沿族参数（Jacobi 常数、振幅等）逐步微调
初始条件、每步用微分修正重新收敛，即生成一族轨道。e2m2e 提供两种
延拓：

\begin{itemize}
  \item \textbf{自然参数延拓}：沿单一参数推进，每步以当前解为初值；
  \item \textbf{伪弧长延拓（PAL）}：增加弧长约束，可越过转向点，
        适合参数--振幅关系非单调的族（如 Halo 的振幅--周期曲线）。
\end{itemize}

失败处理是工程化的：某步不收敛则缩减步长重试，缩到下限仍不收敛则
终止该方向，族中只追加收敛成员，诊断统计总步数 / 成功 / 失败。
L2 NRHO 族采用“折叠后固定 $x_0$ 延拓”的专用编排；Halo 族用固定
$z_0$ 的自然参数延拓，正负 $z_0$ 分北族南族。

\subsection{稳定性分析与不变流形}

周期轨道的稳定性由单值矩阵（一个周期的 STM）决定：
\begin{equation}
  \mat{M} = \stm(T, 0).
\end{equation}
由 Hamiltonian 系统的辛结构，特征值成对出现 $(\lambda, 1/\lambda)$。
Broucke 稳定性指数对每对倒数特征值求和取实部：
\begin{equation}
  \nu = \lambda + \frac{1}{\lambda},
\end{equation}
$|\nu| < 2$ 则该模态稳定，$|\nu| > 2$ 则不稳定且越大越不稳。特征值与
$+1$、$-1$ 及单位圆的关系标识分岔类型（鞍结 / 倍周期 / 环面），支持
族级分岔检测。数值可信度由辛结构残差自校验：
$|\det \mat{M} - 1|$ 与 $\norm{\mat{M}^\top \mat{J} \mat{M} - \mat{J}}$
显著偏大时应收紧积分容差重新分析。

单值矩阵的实特征向量同时给出不变流形方向：稳定流形取 $|\lambda| < 1$
方向、不稳定流形取 $|\lambda| > 1$ 方向。沿轨道相位点用 STM 转运特征
向量、施加 $\pm\varepsilon$ 扰动（典型 50 km/DU）得流形种子，反向 /
正向积分得到流形管，供低能量转移拼接使用
（见第~\ref{sec:transfer}~节）。

\subsection{设计流水线}

把上述算法串起来，一条任务轨道的完整流水线为：

\begin{center}
\begin{tikzpicture}[
  font=\footnotesize,
  stage/.style={draw=e2blue, thick, rounded corners=1.5mm, fill=e2light,
    minimum height=0.85cm, align=center, inner sep=4pt},
  arr/.style={-{Stealth[length=2mm]}, thick, e2blue},
]
\node[stage] (s1) {初猜\\Richardson/种子};
\node[stage, right=0.28cm of s1] (s2) {微分修正\\CR3BP 收敛};
\node[stage, right=0.28cm of s2] (s3) {延拓\\生成轨道族};
\node[stage, right=0.28cm of s3] (s4) {相位定位\\+ 节点采样};
\node[stage, right=0.28cm of s4] (s5) {批量转系\\会合系$\to$J2000};
\node[stage, right=0.28cm of s5] (s6) {星历修正\\多重打靶};
\node[stage, right=0.28cm of s6] (s7) {高保真\\传播};
\draw[arr] (s1) -- (s2);
\draw[arr] (s2) -- (s3);
\draw[arr] (s3) -- (s4);
\draw[arr] (s4) -- (s5);
\draw[arr] (s5) -- (s6);
\draw[arr] (s6) -- (s7);
\end{tikzpicture}
\end{center}

数值重活（传播、打靶、批量坐标转换、批量星历查询）全部在 Rust 侧，
Python 只做编排；四条典型轨道（LLO、ELFO、DRO、Halo）的端到端生成
在 release 构建下总计不到 1 分钟。

\begin{keybox}{NominalOrbit：设计与控制之间的契约}
设计完成的标称轨道以 NominalOrbit 结构交付：等间隔状态表 + Floquet 基
+ 投影因子。轨道保持算法直接消费该结构——设计阶段的 Floquet 分析结果
不需要在控制阶段重算，两个环节通过同一份数据契约解耦。
\end{keybox}
